% ==========================================================================
% Section 1: Introduction — Reframed for methodology + cautionary finding
% ==========================================================================

\section{Introduction}

Benchmark contamination---the presence of test items in training
data---is widely suspected of inflating LLM performance on standard
evaluations~\cite{zhang2024gsm1k,dekoninck2024constat,mmlucf2024}.
The standard remedy is to generate novel benchmark variants that test
the same reasoning while eliminating memorization advantages. But a
critical question has received insufficient attention: \emph{how do we
know the variant answers are correct?}

We present \textbf{Isomorph-Eval}, a framework for generating
structurally equivalent benchmark variants via $\tau$-isomorphism, a
graph-theoretic definition of structural equivalence. In applying
the framework to GSM8K~\cite{cobbe2021gsm8k}, we discovered that the
standard approach to computing variant answers---annotation-chain
replay, in which new values are substituted into the original
solution chain---produces systematically wrong answers in 41\% of
items. These errors created an apparent contamination signal of
$\Delta \approx +0.42$ that was \emph{uniform across all models},
including a comparison model whose training data composition has not
been fully disclosed.
After independent verification and pipeline repair, the signal
collapsed to $\Delta \approx +0.04$--$0.11$.

This finding motivates a reframing: the primary contribution of this
work is not contamination detection, but rather the identification of
eight bug classes in variant generation pipelines that produce false
contamination signals, and the development of verification methods
that prevent them.

\paragraph{Background.}
Prior work on benchmark contamination detection falls into two
categories. \emph{Parallel benchmark} approaches construct entirely
new test sets: GSM1K~\cite{zhang2024gsm1k} reported 8--13\% performance
drops for open-source models on a manually curated GSM8K parallel;
MMLU-CF~\cite{mmlucf2024} found a 14.6-point drop for GPT-4o on a
contamination-free MMLU variant. \emph{Statistical detection} approaches
such as ConStat~\cite{dekoninck2024constat} and
PaCoST use performance gaps on rephrased items as contamination evidence.
A complementary line of work has established that benchmark scores
are also distorted by aggregation methodology: our prior Evaluation
Failure Scaling Law (EFSL)~\cite{kang2026efsl} showed that simple
averaging degrades to $\rho = 0.24$ under the joint effect of data
sparsity and item difficulty heterogeneity, while
PSN-IRT~\cite{zhou2026psnirt} and MedIRT~\cite{luo2025medirt}
demonstrated that IRT-based rankings are more robust.

All of these approaches assume that the alternative evaluation---whether
a parallel benchmark or a rephrased variant---has correct answers. None
systematically verify this assumption. GSM-Symbolic~\cite{mirzadeh2024gsmsymbolic}
comes closest, substituting symbolic values into GSM8K solution
templates. VarBench reports that Llama~3 Instruct~8B drops from
75.8\% on GSM8K to 39.4\% on GSM+ (48.3\% relative drop) under
dynamic variable perturbation~\cite{qian2024varbench}, underscoring
how difficult it is to separate perturbation sensitivity from
contamination.

\paragraph{$\tau$-Isomorphism.}
We formalize the concept of structural equivalence as
\emph{$\tau$-isomorphism} (Section~\ref{sec:isomorphism}): two
benchmark items are $\tau$-isomorphic if there exists a graph
isomorphism between their reasoning graphs that preserves operation
types, dependency structure, and complexity weights. Unlike semantic
similarity heuristics used by prior rephrasing approaches,
$\tau$-isomorphism is a deterministic, verifiable structural condition.
The \emph{Isomorphic Engine} (Section~\ref{sec:engine}) implements a
four-stage pipeline (Parser $\to$ Mutator $\to$ Generator $\to$
Verifier) that automatically produces verified $\tau$-isomorphic
variants with answers computed by forward graph execution.

\paragraph{The cautionary finding.}
Applying the Isomorphic Engine to a stratified sample of 100 GSM8K
items, we found that annotation-chain replay---the mechanism that
computes variant answers by substituting new values into the
original solution annotations---fails in three systematic ways:
\begin{enumerate}[leftmargin=*, nosep]
\item \textbf{Value conflation}: the same numeric value serves
multiple semantic roles in the problem (e.g., ``5 cars'' and
``\$5 each''), but the solution chain merges them into a single
graph node; mutating the node changes both uses in the graph but only
one in the text.
\item \textbf{Phantom leaf values}: the solution chain introduces
intermediate results as leaf nodes (e.g., \texttt{<<40=40>>}) that
do not appear in the question text; mutating these phantom values
produces wrong answers because the graph is structurally incomplete.
\item \textbf{Positional replacement collisions}: sequential text
replacement causes new values to shadow old values of other nodes,
producing text that does not match the graph computation.
\end{enumerate}
These three failure modes, along with five additional bug classes
(magnitude inflation, answer invariance, partial name replacement,
node reuse, and inconsistent entity replacement), produced wrong or
corrupted variants in 41 of 100~items.
All five evaluated models---including a comparison model---appeared
uniformly ``contaminated'' ($\Delta \approx +0.42$) because they were
penalized for \emph{correctly solving} variants whose expected answers
were wrong.

\paragraph{Corrected results.}
After fixing the generation pipeline and independently verifying all
variant answers via forward graph execution and filtering variants
with inconsistent entity replacement, we evaluate five LLMs on
63 entity-clean verified items. Four of five models show no significant performance gap
between originals and variants
($\Delta_{\mathrm{iso}} \in [+0.036, +0.047]$, all $p > 0.05$).
The remaining model, Llama~3.1~8B, shows a mild gap
($\Delta_{\mathrm{iso}} = +0.114$, $p = 0.003$), far smaller than
the Llama~3 Instruct~8B drop reported by VarBench (75.8\% on GSM8K
to 39.4\% on GSM+, 48.3\% relative)~\cite{qian2024varbench}.
GPT-OSS~120B serves as a comparison model: while its training data is
not fully disclosed, its open-weight status and STEM-focused training
make GSM8K inclusion likely, so its small $\Delta_{\mathrm{iso}}$
(+0.036, $p = 0.090$) suggests the $\tau$-isomorphic variants are not
systematically harder rather than confirming absence of contamination.

\paragraph{Contributions.}
\begin{enumerate}[leftmargin=*, nosep]
\item We formalize \emph{$\tau$-isomorphism} as a verifiable
graph-theoretic condition for benchmark item equivalence
(Section~\ref{sec:isomorphism}).

\item We design and release the \emph{Isomorphic Engine}, a
four-stage pipeline that produces multiple verified variants per successfully parsed item with
automatically computed answers and deterministic structural
verification (Section~\ref{sec:engine}).

\item We identify \textbf{eight bug classes} in annotation-chain
replay that produce false contamination signals, and demonstrate that
\emph{independent answer verification via forward graph execution} is
necessary for any benchmark variant generation framework
(Section~\ref{sec:audit}).

\item We provide empirical evidence that on verified data, current
LLMs show robust reasoning transfer to novel numeric contexts, with
remaining gaps consistent with perturbation sensitivity rather than
contamination (Section~\ref{sec:results}).

\item We define $\Delta_{\mathrm{iso}}^{\mathrm{IRT}}$, a
DIF-based metric for measuring perturbation sensitivity with IRT
discrimination weighting (Section~\ref{sec:delta}).
\end{enumerate}

\noindent We release the complete pipeline as open-source
software.\footnote{Code and data:
\url{https://github.com/testofschool/isomorph-eval}}
